\documentclass{beamer}

\usepackage{graphicx}
\usepackage{german}
\usepackage{amsmath}
\usepackage{amssymb}
\usepackage{amsthm}
\usepackage{url}
\usepackage{microtype}

\newcommand{\Z}{\mathbb{Z}}
\newcommand{\R}{\mathbb{R}}
\newcommand{\C}{\mathbb{C}}
\newcommand{\Q}{\mathbb{Q}}
\newcommand{\N}{\mathbb{N}}

\newtheorem{thm}{Satz}

\title{Approximative Darstellung reeller Zahlen}
\date{\today}
\author{Max Recker\\ \url{max.recker@uni-bielefeld.de}}

\begin{document}

%------------------------------------------------
\begin{frame}
\titlepage
\end{frame}

%------------------------------------------------
\begin{frame}{Übersicht}
\begin{enumerate}
\item 4.1 b-adische Darstellung reeller Zahlen
\item 4.2 Maschinenzahlen
\item 4.3 Rundung
\item 4.4 Maschinenzahlarithmetik
\end{enumerate}
\end{frame}

%------------------------------------------------
\begin{frame}{4.1 Motivation}
Ziel des Kapitels:

\bigskip

Jede reelle Zahl soll in einer Basis $b\geq 2$ dargestellt werden.

\bigskip

Bekannte Basen:
\begin{itemize}
\item $b=10$ : Dezimalsystem
\item $b=2$ : Binärsystem
\item $b=16$ : Hexadezimalsystem
\end{itemize}

\bigskip

Diese Darstellung ist die Grundlage für spätere Maschinenzahlen.
\end{frame}

%------------------------------------------------
\begin{frame}{Satz 4.1: b-adische Darstellung}
\small
\begin{thm}
Sei $b\in\N$, $b\geq 2$, und $x\in\R\setminus\{0\}$.

Dann existieren eindeutig bestimmte Zahlen
\[
E\in\Z,\quad \sigma\in\{-1,1\},\quad z_i\in\{0,1,\dots,b-1\},
\]
so dass
\[
x=\sigma\cdot b^E\cdot\left(\sum_{i=0}^{\infty} z_i b^{-i}\right).
\]
\end{thm}

\bigskip

\textbf{Merke:} Jede Zahl besteht aus Vorzeichen, Größe und Ziffernfolge.
\end{frame}

%------------------------------------------------
\begin{frame}{Erklärung der Bestandteile}
\[
x=\sigma\cdot b^E\cdot\left(\sum_{i=0}^{\infty} z_i b^{-i}\right)
\]

\bigskip

\begin{itemize}
\item $\sigma$ = Vorzeichen
\item $b$ = Basis
\item $E$ = Exponent / Größenordnung
\item $z_i$ = $i$-te Ziffer
\item $b^{-i}$ = Stellenwert
\end{itemize}

\bigskip

Also:
\[
\text{Zahl}=\text{Vorzeichen}\cdot\text{Größe}\cdot\text{Mantisse}
\]
\end{frame}

%------------------------------------------------
\begin{frame}{Ziffern im jeweiligen Zahlensystem}
Es gilt immer:
\[
z_i\in\{0,1,\dots,b-1\}
\]

\bigskip

Beispiele:

\begin{itemize}
\item Basis $10$: Ziffern $0$ bis $9$
\item Basis $2$: Ziffern $0,1$
\item Basis $8$: Ziffern $0$ bis $7$
\item Basis $16$: Ziffern $0$ bis $15$
\end{itemize}

\bigskip

Jede Ziffer darf also nur Werte annehmen, die zur Basis passen.
\end{frame}

%------------------------------------------------
\begin{frame}{Bemerkung 4.1: Warum existiert die Summe?}
Wir betrachten:
\[
\sum_{i=0}^{\infty} z_i b^{-i}
\]

\bigskip

Warum existiert der Grenzwert?

\begin{enumerate}
\item Alle Summanden sind nichtnegativ.
\item Die Partialsummen wachsen monoton.
\item Die Summe ist nach oben beschränkt.
\end{enumerate}

Denn:
\[
z_i\leq b-1
\]

also:
\[
\sum_{i=0}^{\infty} z_i b^{-i}
\leq
\sum_{i=0}^{\infty}(b-1)b^{-i}
\]
\end{frame}

%------------------------------------------------
\begin{frame}{Wichtige Folgerung}
Für die Mantisse gilt:
\[
1\leq \sum_{i=0}^{\infty} z_i b^{-i}<b
\]

\bigskip

Begründung:

\begin{itemize}
\item $z_0\neq 0$ liefert die untere Grenze
\item keine Ziffer größer als $b-1$ liefert die obere Grenze
\end{itemize}

\bigskip

\textbf{Wichtig:} Dieser Bereich $[1,b)$ wird später immer wieder gebraucht.
\end{frame}

%------------------------------------------------
\begin{frame}{Beweis zu Satz 4.1 -- Schritt 1}
\textbf{Ziel:} Darstellung von $x$ konstruieren.

\bigskip

Zuerst das Vorzeichen:
\[
\sigma=
\begin{cases}
-1,& x<0\\[0.2cm]
1,& x>0
\end{cases}
\]

\bigskip

Danach genügt es, nur noch $|x|$ zu betrachten.
\end{frame}

%------------------------------------------------
\begin{frame}{Beweis zu Satz 4.1 -- Schritt 2}
Nun bringen wir die Zahl in den Standardbereich $[1,b)$.

\bigskip

Wir wählen:
\[
E:=\lfloor \log_b|x|\rfloor
\]

und setzen:
\[
a_0:=b^{-E}|x|
\]

Dann gilt:
\[
1\leq a_0<b
\]

\bigskip

\textbf{Merke:} $a_0$ heißt normierte Mantisse.
\end{frame}

%------------------------------------------------
%------------------------------------------------
\begin{frame}{Beweis zu Satz 4.1 -- Schritt 3}
\small
\textbf{Satz 4.1:} Für jedes $x\in\R\setminus\{0\}$ existiert eine normalisierte b-adische Darstellung.

\bigskip

\textbf{Ziel von Schritt 3:} Die Ziffern $z_i$ konstruieren.

\bigskip

Wir setzen
\[
a_0:=b^{-E}|x|
\]
und definieren rekursiv für $i\in\N\cup\{0\}$:
\[
z_i:=\lfloor a_i\rfloor,
\qquad
a_{i+1}:=b(a_i-z_i).
\]

\bigskip

\textbf{Idee:}
\begin{enumerate}
\item $z_i$ ist der ganzzahlige Teil von $a_i$, also die nächste Ziffer.
\item $a_i-z_i$ ist der Rest nach dem Komma.
\item Durch Multiplikation mit $b$ wird die nächste Stelle nach vorne geholt.
\end{enumerate}
\end{frame}

%------------------------------------------------
\begin{frame}{Beweis zu Satz 4.1 -- Schritt 4}
\small
\textbf{Satz 4.1:} Für jedes $x\in\R\setminus\{0\}$ existiert eine normalisierte b-adische Darstellung.

\bigskip

Aus
\[
1\le a_0<b
\]
folgt direkt:
\[
z_0\neq 0.
\]

Außerdem gilt für alle $i$:
\[
0\le a_i<b
\quad\Longrightarrow\quad
z_i=\lfloor a_i\rfloor\in\{0,1,\dots,b-1\}.
\]

\bigskip

\textbf{Wichtig:}
\begin{itemize}
\item Die Ziffern liegen genau im erlaubten Bereich.
\item Die erste Ziffer ist nicht null.
\item Damit passt die Konstruktion zur Behauptung des Satzes.
\end{itemize}
\end{frame}

%------------------------------------------------
\begin{frame}{Beispiel zum Konstruktionsverfahren}
\small
\textbf{Satz 4.1:} Für jedes $x\in\R\setminus\{0\}$ existiert eine normalisierte b-adische Darstellung.

\bigskip

\textbf{Beispiel in Basis $10$:}
\[
a_0=3{,}14
\]

Dann:
\[
z_0=\lfloor 3{,}14\rfloor=3
\]
\[
a_1=10(3{,}14-3)=1{,}4
\]
\[
z_1=\lfloor 1{,}4\rfloor=1
\]
\[
a_2=10(1{,}4-1)=4
\]

\bigskip

Also entstehen Schritt für Schritt die Ziffern
\[
3,1,4,\dots
\]

\textbf{Merke:} Genau so liest man die Ziffern der Mantisse nacheinander ab.
\end{frame}

%------------------------------------------------
\begin{frame}{Beweis zu Satz 4.1 -- Schritt 5}
\small
\textbf{Satz 4.1:} Für jedes $x\in\R\setminus\{0\}$ existiert eine normalisierte b-adische Darstellung.

\bigskip

\textbf{Entscheidende Behauptung:} Für alle $n\in\N\cup\{0\}$ gilt
\[
a_0=\sum_{i=0}^{n} z_i b^{-i}+a_{n+1}b^{-n-1}.
\]

\bigskip

\textbf{Bedeutung der Formel:}
\begin{itemize}
\item \(\sum_{i=0}^{n} z_i b^{-i}\) enthält die schon bekannten Ziffern.
\item \(a_{n+1}b^{-n-1}\) ist der Rest.
\item Für großes $n$ wird der Rest immer kleiner.
\end{itemize}

\bigskip

\textbf{Idee:} So nähert sich die endliche Teilsumme immer besser an \(a_0\) an.
\end{frame}

%------------------------------------------------
\begin{frame}{Beweis zu Satz 4.1 -- Induktion: Anfang}
\small
\textbf{Zu zeigen:}
\[
a_0=\sum_{i=0}^{n} z_i b^{-i}+a_{n+1}b^{-n-1}
\quad \text{für alle } n\in\N\cup\{0\}.
\]

\bigskip

\textbf{1. Induktionsanfang:} $n=0$

Dann lautet die Behauptung:
\[
a_0=z_0b^0+a_1b^{-1}.
\]

Da
\[
a_1=b(a_0-z_0)
\]
gilt, folgt
\[
\frac{a_1}{b}=a_0-z_0.
\]

Also:
\[
z_0+\frac{a_1}{b}=z_0+(a_0-z_0)=a_0.
\]

Damit stimmt die Formel für $n=0$.
\end{frame}

%------------------------------------------------
\begin{frame}{Beweis zu Satz 4.1 -- Induktion: Schritt}
\small
\textbf{Induktionsannahme:}
\[
a_0=\sum_{i=0}^{n-1} z_i b^{-i}+a_n b^{-n}.
\]

\bigskip

\textbf{2. Induktionsschritt:} Zeige die Formel für $n$.

Aus der Definition
\[
a_{n+1}=b(a_n-z_n)
\]
folgt
\[
a_n=z_n+\frac{a_{n+1}}{b}.
\]

Einsetzen in die Induktionsannahme liefert:
\[
a_0=\sum_{i=0}^{n-1} z_i b^{-i}+\left(z_n+\frac{a_{n+1}}{b}\right)b^{-n}.
\]

Ausmultiplizieren:
\[
a_0=\sum_{i=0}^{n-1} z_i b^{-i}+z_n b^{-n}+a_{n+1}b^{-n-1}.
\]

Also:
\[
a_0=\sum_{i=0}^{n} z_i b^{-i}+a_{n+1}b^{-n-1}.
\]

Damit ist die Behauptung bewiesen.
\end{frame}

%------------------------------------------------
\begin{frame}{Beweis zu Satz 4.1 -- Schritt 6}
\small
\textbf{Wichtige Behauptung:}
\[
a_0=\sum_{i=0}^{n} z_i b^{-i}+a_{n+1}b^{-n-1}.
\]

\bigskip

Wir lassen jetzt $n\to\infty$ gehen.

Dann verschwindet der Restterm
\[
a_{n+1}b^{-n-1}\to 0,
\]
denn:
\[
0\le a_{n+1}<b
\]
und daher
\[
0\le a_{n+1}b^{-n-1}<b\cdot b^{-n-1}=b^{-n}\to 0.
\]

Also folgt
\[
a_0=\sum_{i=0}^{\infty} z_i b^{-i}.
\]

Mit
\[
x=\sigma\cdot b^E\cdot a_0
\]
erhalten wir schließlich
\[
x=\sigma\cdot b^E\cdot \left(\sum_{i=0}^{\infty} z_i b^{-i}\right).
\]

\textbf{Damit ist die Darstellung konstruiert.}
\end{frame}

%------------------------------------------------
\begin{frame}{Beweis zu Satz 4.1 -- Schritt 7: Existenz}
\small
\textbf{Noch zu zeigen:}
\[
\{i\in\N\mid z_i\neq b-1\}\ \text{ist unendlich.}
\]

\bigskip

Angenommen, das wäre falsch.

Dann gäbe es ein \(n_0\in\N\), so dass für alle \(i>n_0\) gilt:
\[
z_i=b-1.
\]

Dann wäre
\[
a_0=\sum_{i=0}^{n_0} z_i b^{-i}+\sum_{i=n_0+1}^{\infty}(b-1)b^{-i}.
\]

Die zweite Summe kann man ausrechnen zu
\[
b^{-n_0}.
\]

Also:
\[
a_0=\sum_{i=0}^{n_0} z_i b^{-i}+b^{-n_0}.
\]
\end{frame}

%------------------------------------------------
\begin{frame}{Beweis zu Satz 4.1 -- Schritt 8: Widerspruch}
\small
\textbf{Annahme zum Widerspruch:} Ab einem Index gilt immer \(z_i=b-1\).

\bigskip

Aus der vorher bewiesenen Behauptung folgt aber auch:
\[
a_0=\sum_{i=0}^{n_0} z_i b^{-i}+a_{n_0+1}b^{-n_0-1}.
\]

Vergleich mit
\[
a_0=\sum_{i=0}^{n_0} z_i b^{-i}+b^{-n_0}
\]
liefert:
\[
a_{n_0+1}b^{-n_0-1}=b^{-n_0}.
\]

Also:
\[
a_{n_0+1}=b.
\]

Das ist unmöglich, denn wir wissen bereits:
\[
0\le a_i<b.
\]

\bigskip

\textbf{Widerspruch.} Also ist
\[
\{i\in\N\mid z_i\neq b-1\}
\]
unendlich.

Damit ist die \textbf{Existenz} vollständig bewiesen.
\end{frame}

%------------------------------------------------
\begin{frame}{Beweis zu Satz 4.1 -- Schritt 9: Eindeutigkeit}
\small
\textbf{Nun zur Eindeutigkeit.}

Angenommen, es gäbe zwei Darstellungen:
\[
x=\sigma\cdot b^E\cdot \sum_{i=0}^{\infty} y_i b^{-i}
=
\sigma\cdot b^E\cdot \sum_{i=0}^{\infty} z_i b^{-i}.
\]

Sei $n$ der kleinste Index mit
\[
y_n\neq z_n.
\]

Ohne Einschränkung gelte:
\[
y_n+1\le z_n.
\]

Dann ist die zweite Darstellung ab der Stelle $n$ mindestens „größer“ als die erste.

\bigskip

\textbf{Idee:} Daraus folgt am Ende
\[
\sum_{i=0}^{\infty} y_i b^{-i}<\sum_{i=0}^{\infty} z_i b^{-i},
\]
also ein Widerspruch.
\end{frame}

%------------------------------------------------
\begin{frame}{Beweis zu Satz 4.1 -- Schritt 10: Eindeutigkeit}
\small
Aus der Annahme \(y_n+1\le z_n\) folgt:
\[
\sum_{i=0}^{\infty} y_i b^{-i}
\le
\sum_{i=0}^{n-1} y_i b^{-i}+(z_n-1)b^{-n}+\sum_{i=n+1}^{\infty} y_i b^{-i}.
\]

Da alle Ziffern höchstens \(b-1\) sind, gilt weiter:
\[
\sum_{i=n+1}^{\infty} y_i b^{-i}
<
\sum_{i=n+1}^{\infty}(b-1)b^{-i}.
\]

Zusammen ergibt sich:
\[
\sum_{i=0}^{\infty} y_i b^{-i}
<
\sum_{i=0}^{n-1} z_i b^{-i}+z_n b^{-n}
\le
\sum_{i=0}^{\infty} z_i b^{-i}.
\]

Also wären die beiden Darstellungen verschieden.

Das widerspricht aber der Annahme, dass beide dieselbe Zahl \(x\) darstellen.

\bigskip

\textbf{Damit ist auch die Eindeutigkeit bewiesen.}
\end{frame}

%------------------------------------------------
\begin{frame}{Beispiel 4.2}
\small
\textbf{Beispiel aus dem Skript für \(b=2\):}

\[
\frac13
=
2^{-2}\left(2^0+2^{-2}+2^{-4}+2^{-6}+\dots\right)
=
(0,01)_2
\]

\bigskip

Außerdem:
\[
8{,}1
=
2^3\left(2^0+2^{-7}+2^{-8}+2^{-11}+2^{-12}+\dots\right)
=
(1000,00011\dots)_2
\]

\bigskip

\textbf{Wichtig:}
Auch bekannte Dezimalzahlen können in Basis 2 eine unendliche Darstellung besitzen.
\end{frame}

%------------------------------------------------
\begin{frame}{Zusammenfassung zu Satz 4.1}
\small
\textbf{Satz 4.1:} Jede reelle Zahl \(x\neq 0\) besitzt eine eindeutig bestimmte normalisierte b-adische Darstellung.

\bigskip

\textbf{Beweisstruktur:}
\begin{enumerate}
\item Vorzeichen bestimmen
\item Zahl normieren
\item Ziffern rekursiv konstruieren
\item zentrale Behauptung per Induktion beweisen
\item Existenz vollständig zeigen
\item Eindeutigkeit mit Widerspruch zeigen
\end{enumerate}

\bigskip

\textbf{Bedeutung für das nächste Kapitel:}
Diese Darstellung ist die Grundlage der Maschinenzahlen.
\end{frame}
